\section{Introduction}\label{sec:introduction}

A polynomial difference theorem asks how large a set of integers can be
if the differences of its elements avoid a prescribed polynomial
sequence. Restricting the polynomial's argument to primes imposes an
additional arithmetic condition: a root modulo a given modulus must
be compatible with a prime residue class. We prove a power-saving bound
when these local conditions hold at every modulus.

Write $\mathbb P$ for the positive rational primes and
$[N]=\{1,\ldots,N\}$. An integer polynomial $h$ is
\emph{prime-intersective} if, for every positive integer $q$, there is
an integer $r$ with
\[
 (r,q)=1,\qquad h(r)\equiv0\pmod q.
\]
For a set $A$ of integers, write $A-A=\{a-b:a,b\in A\}$.

\begin{theorem}\label{thm:main}
Let $h\in\Z[x]$ be prime-intersective, of degree at least two and with
positive leading coefficient. There exist constants $c_h>0$ and
$C_h\ge1$ such that, for every positive integer $N$ and every
$A\subseteq[N]$,
\[
 (A-A)\cap\{h(p):p\in\mathbb P\}\subseteq\{0\}
 \quad\Longrightarrow\quad
 |A|\le C_h N^{1-c_h}.
\]
\end{theorem}

The convention at zero allows a prime root of $h$: a value $h(p)=0$
places no restriction on $A$. All constants in the theorem depend only
on the fixed polynomial. Increasing $C_h$ incorporates the finitely
many small values of $N$.

\subsection{History and main ideas}

The starting point is the Furstenberg--S\'ark\"ozy theorem: every subset
of the integers with positive upper density contains two elements whose
difference is a nonzero square. Furstenberg's ergodic proof and
S\'ark\"ozy's quantitative Fourier argument established this independently
\cite{Furstenberg1977,Sarkozy1978}. For a nonconstant integer polynomial,
the unavoidable local condition is a root modulo every positive integer.
Kamae and Mend\`es France's difference theorem implies that this
condition also suffices for density decay
\cite{KamaeMendesFrance1978}. Lucier made the general polynomial theorem
quantitative by organizing the iteration around compatible $p$-adic
roots and a family of auxiliary polynomials \cite{Lucier2006}.

Much subsequent work sharpened the density bounds. For square
differences, the double iteration of Pintz, Steiger, and Szemer\'edi
\cite{PintzSteigerSzemeredi1988} was improved by Bloom and Maynard
\cite{BloomMaynard2022}, and Green and Sawhney obtained
$N\exp(-c\sqrt{\log N})$ using an arithmetic level-$d$ inequality
\cite{GreenSawhney2025}. Rice \cite{Rice2019} and Arala
\cite{Arala2024} extended the earlier bounds to general intersective
polynomials. Adajar and collaborators subsequently obtained
$N\exp(-c_{h,\mu}(\log N)^\mu)$ for every $0<\mu<1/2$
\cite{AdajarEtAl2026}. These results concern unrestricted integer
arguments. Restricting the argument to primes makes avoidance a weaker
condition and introduces the unit-root requirement in our definition.

The prime problem already has a substantial history.
S\'ark\"ozy treated the shifts $p-1$ and $p+1$
\cite{Sarkozy1978III}. After successive improvements by Lucier,
Ruzsa--Sanders, and Wang
\cite{Lucier2008,RuzsaSanders2008,Wang2020}, Green proved a power saving
for $p-1$, noting that his argument also applies to $p+1$
\cite[Theorem~1.1 and the following remarks]{Green2024}.
Thorner and Zaman made an admissible exponent explicit
\cite{ThornerZaman2024}.
For nonlinear polynomials, Li and Pan proved a quantitative result when
$h(1)=0$ \cite{LiPan2009}. Rice's theorem covers every
prime-intersective polynomial of degree $k\ge2$, with
\[
 |A|\ll_{h,c}N(\log N)^{-c}
 \qquad\text{for every }c<\frac{1}{2k-2},
\]
and gives a stronger bound for quadratics
\cite[Theorems~1 and~2]{Rice2013}. Rice also records Wierdl's earlier
qualitative characterization by the unit-root condition
\cite[Section~1.2]{Rice2013}. The auxiliary family used below, including
its prime-argument restriction, already appears in
\cite[Section~2.1]{Rice2013}.

Our direct methodological predecessor is the companion
\emph{A power saving for intersective polynomial differences with an exponent depending only on the degree}
\cite{IntersectivePower}. We adapt its reflection-positive tuple laws,
Fourier lifts, signed kernels, and contracting induction to prime
arguments. In particular, the Fourier-lift moment estimate is reproduced
with its proof below. The new requirements concern the compatibility
of the tuple laws with unit residue classes and the realization of
their pair correlations by a signed kernel supported on values at
prime arguments. The example $h(x)=(x-1)^2$ gives a power saving when
only the differences $(p-1)^2$ are forbidden; it therefore implies a
power saving for sets avoiding all nonzero square differences as well.

Here is the mechanism. For each sufficiently large prime we construct
a probability law on a fixed finite tuple of residue labels. A
designated cycle of tuple positions has increments satisfying the
congruence conditions imposed by prime arguments, apart from a component
of very small mass. Reflection positivity means that the joint distribution across
each prescribed cut of the tuple gives a positive semidefinite matrix.
The cuts are chosen so that successive reflections can make all inputs
agree. Repeated Cauchy--Schwarz inequalities then give a multilinear
H\"older inequality and a seminorm. This is the classical chessboard
method of reflection positivity
\cite[Theorem~4.1]{FILS1978}. For related H\"older inequalities and
folding arguments in graph-norm theory, see
\cite{Hatami2010,ConlonLee2017}. We prove the form needed for our
tuple laws directly.

We next lift an interval indicator to the product of these residue
spaces by retaining its rational Fourier coefficients with bounded
denominator. Averaging a product of copies of this lift over the tuple
laws gives a nonnegative quantity that dominates a fixed power of the
set's density. To exploit avoidance, a signed weight on prime polynomial
values is compared with a model kernel encoding the correlations of a
cycle edge; their Fourier transforms are uniformly close. Major arcs match the local
unit distributions; minor arcs are controlled by Vaughan's identity
and Weyl differencing \cite{Vaughan1980,Vaughan1997,Weyl1916}.
Avoidance annihilates the prime-supported kernel on ordered pairs from the set.
The resulting cancellation, together with the seminorm estimate for
restriction to progressions, produces a contraction of the lifted
quantity on shorter intervals. Induction turns that contraction into
the claimed power saving.

The prime-distribution estimate used in the major arcs is derived from
\cite[Theorem~1.1]{ZeroFree}: every Dirichlet $L$-function, uniformly
over all characters and moduli, has no zero in $\Re s>7/8$, with the
principal pole at $s=1$ allowed. This fixed half-plane is a theorem input
from that companion, stronger than the classical zero-free regions
near $\Re s=1$. We prove here its required consequence for primes in
arithmetic progressions and all subsequent steps of the argument.

\subsection{Notation and organization}

Throughout the proof, $h$ is fixed, its degree is $k\ge2$, and its
leading coefficient is $a>0$. Put $e(x)=\exp(2\pi i x)$.
For $\xi\in\Q/\Z$, let $q(\xi)$ be its reduced positive denominator,
with $q(0)=1$. Fourier coefficients use the negative sign in the
exponential. Averages and norms on finite sets use uniform probability
measure unless another measure is specified. Empty products have
value one, and an empty product of spaces is a singleton.

Implicit constants and lower thresholds may depend on $h$ and on the
structural parameters fixed in terms of $h$, but are uniform in the
auxiliary parameter and residue classes introduced below. The notation
$N^{o(1)}$ denotes a quantity bounded by $C_\epsilon N^\epsilon$ for
every $\epsilon>0$, with this same uniformity.

Section~\ref{sec:local} constructs the auxiliary polynomials and their
admissible local values. Section~\ref{sec:tuples} builds the
reflection-positive tuple laws, and Section~\ref{sec:lifts} develops
their multilinear estimates and proves the Fourier-lift moment bound.
Section~\ref{sec:primes} obtains the prime estimates from the stated
zero-free input and elementary differencing. Section~\ref{sec:kernel}
matches the tuple multipliers with a kernel supported on prime
polynomial values. Section~\ref{sec:descent} uses that comparison to
contract a multilinear density quantity and proves
Theorem~\ref{thm:main}.
